\chapter{Introduction and linear algebra}
\label{chap:lecture01}

Robotics brings together mathematical models of motion, measurements from
sensors, and algorithms that determine how a robot should act. Describing a
robot's position, relating measurements to its surroundings, and planning its
motion all require a precise mathematical language. Linear algebra provides
many of the tools used to express these relationships.

This chapter introduces matrices and the operations needed to work with them:
addition, multiplication, transposition, and inversion. It then examines
orthogonal matrices, whose inverses can be obtained simply by transposing.
The examples develop two habits that will remain essential throughout the
study of robotics: checking dimensions before performing an operation and
preserving the order of factors in a matrix product.
\section{Matrices and dimensions}

A \emph{matrix} is an array of numbers arranged in rows and columns. Its
dimensions are written as
\[
  \text{number of rows}\times\text{number of columns}.
\]
Consider the matrix
\begin{equation}
 A=\begin{bmatrix}1&2\\3&1\\0&1\end{bmatrix}\in\R^{3\times2}.
 \label{eq:l01-A}
\end{equation}
It has three rows and two columns. In general, $a_{ij}$ denotes the entry in
row $i$, column $j$; for example, $a_{21}=3$ and $a_{32}=1$.

\section{Matrix addition}

Two matrices can be added only if they have the same dimensions. Their entries
are then added in corresponding positions:
\[
 (A+C)_{ij}=a_{ij}+c_{ij}.
\]
Consider two additional matrices,
\[
 B=\begin{bmatrix}1&2\\1&1\end{bmatrix}\in\R^{2\times2},
 \qquad
 C=\begin{bmatrix}0&1\\0&0\\1&0\end{bmatrix}\in\R^{3\times2}.
\]
The sum $A+B$ is undefined because the dimensions differ. The sum $A+C$ is
defined, and its calculation is
\begin{equation}
 A+C
 =\begin{bmatrix}1+0&2+1\\3+0&1+0\\0+1&1+0\end{bmatrix}
 =\begin{bmatrix}1&3\\3&1\\1&1\end{bmatrix}.
\end{equation}
Addition is commutative when defined: $A+C=C+A$.

\section{Matrix multiplication}

\subsection{Dimension rule and row-by-column calculation}
For matrices $M\in\R^{m\times n}$ and $N\in\R^{n\times p}$, the product
$MN$ is defined and has dimensions $m\times p$:
\[
 \underbrace{(m\times n)(n\times p)}_{\text{inner dimensions must match}}
 \quad\longrightarrow\quad m\times p.
\]
The entry in row $i$, column $j$ is the sum of the products of corresponding
entries of row $i$ of $M$ and column $j$ of $N$:
\begin{equation}
 (MN)_{ij}=\sum_{k=1}^{n}m_{ik}n_{kj}.
\end{equation}
Matrix multiplication is therefore a row-by-column operation, rather than
entrywise multiplication.

For the matrices $A$ and $B$ above, the dimensions are
$(3\times2)(2\times2)$, so $AB$ is a $3\times2$ matrix:
\begin{align}
 AB
 &=\begin{bmatrix}1&2\\3&1\\0&1\end{bmatrix}
   \begin{bmatrix}1&2\\1&1\end{bmatrix}\\
 &=\begin{bmatrix}
  1\cdot1+2\cdot1 & 1\cdot2+2\cdot1\\
  3\cdot1+1\cdot1 & 3\cdot2+1\cdot1\\
  0\cdot1+1\cdot1 & 0\cdot2+1\cdot1
 \end{bmatrix}
 =\begin{bmatrix}3&4\\4&7\\1&1\end{bmatrix}.
 \label{eq:l01-AB}
\end{align}

\subsection{The order of the factors matters}
Reversing the factors gives the dimensions $(2\times2)(3\times2)$. The inner
dimensions do not match, so $BA$ is undefined, even though $AB$ exists.

Even when both orders are defined, the results need not agree. With
\[
 E=\begin{bmatrix}1&2\\3&4\end{bmatrix},
\]
direct multiplication gives
\[
 BE=\begin{bmatrix}7&10\\4&6\end{bmatrix},
 \qquad
 EB=\begin{bmatrix}3&4\\7&10\end{bmatrix}.
\]
Thus $BE\ne EB$. Matrix multiplication is not commutative in general. Some
particular matrices do commute, but factors cannot be exchanged without a
reason that applies to the matrices in question.

\section{The transpose}

Transposing a matrix interchanges its rows and columns. In entry notation,
$(M\trans)_{ij}=m_{ji}$. If $M$ has dimensions $m\times n$, then $M\trans$
has dimensions $n\times m$. For the running example,
\begin{equation}
 A\trans=\begin{bmatrix}1&3&0\\2&1&1\end{bmatrix}\in\R^{2\times3}.
\end{equation}

The transpose of a product reverses the order of its factors:
\begin{equation}
 \boxed{(AB)\trans=B\trans A\trans.}
 \label{eq:l01-transpose-product}
\end{equation}
The dimensions provide a quick check. Here $B\trans A\trans$ has dimensions
$(2\times2)(2\times3)$, producing a $2\times3$ matrix as required. The
incorrect expression $A\trans B\trans$ would have dimensions
$(2\times3)(2\times2)$ and is undefined.

Using the matrices defined above, this identity can be verified directly:
\[
 B\trans A\trans
 =\begin{bmatrix}1&1\\2&1\end{bmatrix}
  \begin{bmatrix}1&3&0\\2&1&1\end{bmatrix}
 =\begin{bmatrix}3&4&1\\4&7&1\end{bmatrix}
 =(AB)\trans.
\]
Dimension checks can rule out an invalid expression, but matching dimensions
alone do not prove that an identity is correct.

\section{Identity matrices and inverses}

The identity matrix $I_n$ has ones on its main diagonal and zeros elsewhere.
For a square matrix $M\in\R^{n\times n}$, an inverse $M^{-1}$ satisfies
\begin{equation}
 MM^{-1}=M^{-1}M=I_n.
\end{equation}
This ordinary, two-sided inverse is defined only for square matrices, and not
every square matrix has one. The condition for invertibility is
\begin{equation}
 \boxed{\det(M)\ne0.}
\end{equation}
A square matrix with zero determinant is called \emph{singular}.

\subsection{A one-by-one example}
For $H=\begin{bmatrix}5\end{bmatrix}$, the inverse is
$H^{-1}=\begin{bmatrix}1/5\end{bmatrix}$, since
\[
 HH^{-1}=H^{-1}H=\begin{bmatrix}1\end{bmatrix}=I_1.
\]
This is the familiar reciprocal operation for a nonzero scalar.

\subsection{A diagonal example and its determinant}
The second example is
\[
 G=\begin{bmatrix}1&0\\0&2\end{bmatrix},
 \qquad
 G^{-1}=\begin{bmatrix}1&0\\0&1/2\end{bmatrix}.
\]
Multiplying verifies the inverse:
\[
 GG^{-1}
 =\begin{bmatrix}
 1\cdot1+0\cdot0&1\cdot0+0\cdot\tfrac12\\
 0\cdot1+2\cdot0&0\cdot0+2\cdot\tfrac12
 \end{bmatrix}
 =\begin{bmatrix}1&0\\0&1\end{bmatrix}=I_2.
\]
The reverse product also equals $I_2$. For a general $2\times2$ matrix,
\[
 \det\begin{bmatrix}a&b\\c&d\end{bmatrix}=ad-bc.
\]
Hence $\det(G)=1\cdot2-0\cdot0=2\ne0$, as required. Reciprocating the
diagonal entries works here because $G$ is diagonal with nonzero diagonal
entries; matrix inversion is not generally an entrywise reciprocal operation.

\section{Orthogonal matrices}

A real square matrix $L$ is \emph{orthogonal} if its inverse equals its
transpose:
\begin{equation}
 \boxed{L^{-1}=L\trans.}
\end{equation}
Equivalently, $LL\trans=L\trans L=I$. This gives a practical way to check
orthogonality: form the transpose, multiply, and verify the identity matrix.
For an orthogonal matrix, an inverse can be obtained by transposing instead of
performing a general inverse calculation.

\subsection{Verifying orthogonality}
Consider
\[
 L=\begin{bmatrix}
 \tfrac12&\tfrac{\sqrt3}{2}\\
 -\tfrac{\sqrt3}{2}&\tfrac12
 \end{bmatrix},
 \qquad
 L\trans=\begin{bmatrix}
 \tfrac12&-\tfrac{\sqrt3}{2}\\
 \tfrac{\sqrt3}{2}&\tfrac12
 \end{bmatrix}.
\]
Their product is
\begin{align}
 LL\trans
 &=\begin{bmatrix}
 \tfrac14+\tfrac34&-\tfrac{\sqrt3}{4}+\tfrac{\sqrt3}{4}\\
 -\tfrac{\sqrt3}{4}+\tfrac{\sqrt3}{4}&\tfrac34+\tfrac14
 \end{bmatrix}\\
 &=\begin{bmatrix}1&0\\0&1\end{bmatrix}=I_2.
\end{align}
Thus $L$ is orthogonal and the displayed $L\trans$ is its inverse.

\subsection{The trigonometric pattern}
The preceding example belongs to the family
\begin{equation}
 Q(\theta)=\begin{bmatrix}
 \sin\theta&\cos\theta\\
 -\cos\theta&\sin\theta
 \end{bmatrix}.
 \label{eq:l01-Q}
\end{equation}
Multiplying by the transpose gives
\begin{align*}
 Q(\theta)Q(\theta)\trans
 &=\begin{bmatrix}
 \sin^2\theta+\cos^2\theta & -\sin\theta\cos\theta+\cos\theta\sin\theta\\
 -\cos\theta\sin\theta+\sin\theta\cos\theta & \cos^2\theta+\sin^2\theta
 \end{bmatrix}\\
 &=I_2.
\end{align*}
The identity $\sin^2\theta+\cos^2\theta=1$ makes this family orthogonal for
every real $\theta$.

The numerical example corresponds to $L=Q(\pi/6)$, since
$\sin(\pi/6)=1/2$ and $\cos(\pi/6)=\sqrt3/2$. An equivalent
parameterization uses cosine on the diagonal:
\[
 P(\varphi)=\begin{bmatrix}
 \cos\varphi&\sin\varphi\\
 -\sin\varphi&\cos\varphi
 \end{bmatrix}.
\]
The relation $Q(\theta)=P(\pi/2-\theta)$ follows from the complementary-angle
identities. Thus the same numerical matrix is $L=P(\pi/3)$. In either
parameterization, orthogonality follows from the same trigonometric identity.
Such matrices provide a foundation for describing changes of orientation
in robotics.
\section{Key identities}
For matrices of compatible dimensions, the operations developed in this
chapter can be collected as follows:
\begin{center}
\begin{tabular}{ll}
\toprule
Operation or property & Rule \\
\midrule
Addition & $(M+N)_{ij}=m_{ij}+n_{ij}$ \\
Multiplication & $(MN)_{ij}=\sum_{k=1}^{n}m_{ik}n_{kj}$ \\
Transpose & $(M\trans)_{ij}=m_{ji}$ \\
Transpose of a product & $(MN)\trans=N\trans M\trans$ \\
Inverse & $MM^{-1}=M^{-1}M=I$ \\
Invertibility of a square matrix & $\det(M)\ne0$ \\
Orthogonality & $M^{-1}=M\trans$ \\
\bottomrule
\end{tabular}
\end{center}
In the multiplication formula, $n$ is the number of columns of $M$ and the
number of rows of $N$. Addition requires equal dimensions, whereas a product
requires only that these inner dimensions match. Inverse and orthogonality
conditions apply to square matrices. These distinctions determine which
algebraic manipulations are valid before any numerical calculation begins.

\clearpage
\section{Additional exercises}
Show your calculations and justify each conclusion. Check dimensions before
performing an operation. These exercises can be completed by hand.

\begin{enumerate}
\item \textbf{Dimensions and valid operations.}
Let $A\in\R^{3\times2}$, $B\in\R^{2\times3}$, and
$C\in\R^{3\times3}$. Determine which expressions are defined and give the
dimensions of each valid result:
\[
 A+B,\qquad AB,\qquad BA,\qquad CA,\qquad AC,\qquad
 A\trans+C,\qquad B+A\trans.
\]

\item \textbf{Addition and multiplication.}
Consider
\[
 M=\begin{bmatrix}2&-1\\0&3\end{bmatrix},\qquad
 N=\begin{bmatrix}1&2\\-1&0\end{bmatrix},\qquad
 v=\begin{bmatrix}1\\-2\end{bmatrix}.
\]
Compute $M+N$, $MN$, $NM$, and $Mv$. Do $M$ and $N$ commute?
Verify numerically that $(M+N)v=Mv+Nv$.

\item \textbf{Transposing a product.}
Let
\[
 U=\begin{bmatrix}1&0&-1\\2&1&3\end{bmatrix},\qquad
 V=\begin{bmatrix}1&2\\0&-1\\2&0\end{bmatrix}.
\]
Compute $UV$ and $(UV)\trans$, then independently compute
$V\trans U\trans$ and compare the results. Is $U\trans V\trans$
defined? Explain why it cannot equal $(UV)\trans$ in this example.

\item \textbf{Determinants and inverses.}
For
\[
 D=\begin{bmatrix}4&0\\0&-2\end{bmatrix},\qquad
 F=\begin{bmatrix}1&2\\2&4\end{bmatrix},\qquad
 K=\begin{bmatrix}2&1\\1&1\end{bmatrix},
\]
use determinants to decide which matrices are invertible. Find $D^{-1}$.
Verify that $\begin{bmatrix}1&-1\\-1&2\end{bmatrix}$ is the inverse of
$K$ by multiplying in both orders.

\item \textbf{Testing orthogonality.}
Consider
\[
 R=\begin{bmatrix}\tfrac35&-\tfrac45\\[2pt]\tfrac45&\tfrac35\end{bmatrix},
 \qquad S=\begin{bmatrix}1&0\\0&-1\end{bmatrix},
 \qquad T=\begin{bmatrix}1&1\\0&1\end{bmatrix}.
\]
Test each matrix for orthogonality using its transpose. For each orthogonal
matrix, write its inverse. Is an invertible matrix necessarily orthogonal?
Support your answer with one of these examples.

\item \textbf{A product of orthogonal matrices.}
Suppose $P$ and $Q$ are real orthogonal $n\times n$ matrices. Use the
transpose-of-a-product rule to show that $(PQ)(PQ)\trans=I_n$.
Deduce that $PQ$ is orthogonal and express $(PQ)^{-1}$ in terms of
$P\trans$ and $Q\trans$, paying attention to their order.
\end{enumerate}
